\scanpage{112}
\deleted{a) Si $F$ spécial de \unsure{pds|poids}}

\uline{Proposition} 3.3 \quad a) Stabilité de \textquotedblleft eff. adm. de poids $\subset P$\textquotedblright{} par sous-faisceau constructible, quotient constructible, \inserted{sommes finies,} extensions.

b) Si $F,F'$ sont eff. adm à poids $\subset P,P'$, alors $F\otimes F'$ est admissible à poids $\subset P+P'=\{i+i'\mid i\in P,i'\in P'\}$.

c) Stabilité par chgt de base

d) Si $F$ est \deleted{spécial} eff$^{\mathrm t}$ admissible \uline{constant} \uline{tordu} à poids $\subset P$, \deleted{\unsure{et}} si $i\in P\Rightarrow i\le n$, alors $\check F(\unsure{-n|n})$ est à \inserted{eff$^{\mathrm t}$} admissible (constant tordu) à poids $\subset 2n-P=\{n-i\mid i\in P\}$.
\margin{d) modulo Lefschetz}

\uline{Dém} \quad a) b) c) sont immédiats, on utilise les assertions correspondantes de 2.1. Notons d'ailleurs :

% Left of 3.2; the arrowhead is beside the top line of the sheet.
\figure{\begin{tikzpicture}[x=.42cm,y=.33cm]
\draw[->] (1,0) .. controls (1.5,2.8) and (-2,6.4) .. (-1,8.3) .. controls (-.7,9.0) and (.2,9.6) .. (.7,9.8);
\draw (1.8,.5) .. controls (.8,.4) and (.9,-.7) .. (.9,-1.1) -- (.9,-3.0) .. controls (.9,-3.6) and (1.2,-3.6) .. (1.8,-3.45);
\end{tikzpicture}}
3.2. \inserted{Supposons $S$ qu. cpt qu. sép.} Pour que $F$ soit eff$^{\mathrm t}$ adm. à poids $\subset P$, il f et s. que l'on puisse trouver une partition de $S$ en \deleted{\unsure{fermés}} sous-schémas $S_i\overset{\alpha_i}{\lhook\joinrel\longrightarrow} S$, $\alpha_i$ imm. de prés. finie, avec \deleted{\unsure{des}} $F|S_i=F_i$ admettant une filtration dont les quotients successifs sont spéciaux de poids $\in P$.

Comme, pour $F$ constant tordu, la formation de $\check F$ commute au changement de base, \inserted{3.\unsure{2|?}} d) résulte de 2.2.

4.1 Un faisceau $F$ est dit eff$^{\mathrm t}$ admissible \deleted{\unsure{poids}} de niveaux $\le\unsure{\nu|v}$ si \deleted{il est effectivement admissible, et} dans 3.1, on peut prendre $G$ \deleted{\unsure{de poids $j\in$}} de la forme $G'(-k)$, $k\ge0$ où $G'$ est spécial de poids \deleted{\unsure{$j$}} $\le\nu$, \deleted{\unsure{on|---}} On dit que $F$ est eff$^{\mathrm t}$ admissible \textquotedblleft de poids $\subset P$, et de niveau $\le\nu$\textquotedblright{} si on peut choisir dans 3.1. $G'$ \deleted{\unsure{à la}} spécial de poids $j\le\nu$, avec \deleted{\unsure{poids $G$}} $j+2k\in P$. \uline{Attention}, à priori cela est plus fort que la \uline{conjonction} de \textquotedblleft de poids $\subset P$\textquotedblright{} et \textquotedblleft de niveau $\le\nu$\textquotedblright. \deleted{\unsure{Pour cette notion conjointe, on a un énoncé :} \unsure{4.2 a) b) c) d)|4.2 a) b) c)}} \inserted{On notera que si \deleted{\unsure{$P\le\nu$}}, $\nu$ est} un entier tel que $P\le\nu$ (i.e $j\in P\Rightarrow j\le\nu$), alors \textquotedblleft eff$^{\mathrm t}$ adm. de poids $\in P$\textquotedblright{} $\to$ \textquotedblleft id et à niveau $\le\nu$\textquotedblright

% slanted
\margin{bloquer après 3.2, avec \unsure{énoncé analogue} : 3.2, pour faire sortir 3.3 en $=$ temps que 4.2.}
% archive sheet number: 111
